\documentclass[11pt,a4paper]{article}
\usepackage[utf8]{inputenc}
\usepackage[T1]{fontenc}
\usepackage[french]{babel}
\usepackage{amsmath,amssymb}
\usepackage[top=2cm,bottom=2.5cm,left=2.5cm,right=2cm]{geometry}
\usepackage{enumitem}
\usepackage{booktabs}

\begin{document}

\begin{center}
{\large\bfseries Préparation au concours d'entrée à l'IFORD -- Yaoundé}\\[2mm]
{\Large\bfseries Concours TYPE A -- Épreuve de PROBABILITÉS ET STATISTIQUE}\\[2mm]
Sujet d'entraînement n\textsuperscript{o}\,1 \quad --\quad Durée : 4 heures\\[1mm]
\emph{Sujet rédigé pour l'entraînement, conforme au programme officiel. Ce n'est pas un sujet officiel.}
\end{center}
\hrule\medskip

\noindent\textbf{Consignes.} Les exercices sont indépendants. Calculatrice non programmable autorisée.
Donner les résultats décimaux à $10^{-2}$ près.

\section*{Exercice 1 -- Statistique descriptive (6 points)}
Une enquête auprès de $100$ femmes a relevé l'âge à la première naissance :
\begin{center}
\begin{tabular}{lccccc}
\toprule
Âge (ans) & $[15;20[$ & $[20;25[$ & $[25;30[$ & $[30;35[$ & $[35;40[$\\
\midrule
Effectif & 12 & 28 & 35 & 18 & 7\\
\bottomrule
\end{tabular}
\end{center}
\begin{enumerate}
  \item Préciser la population, l'unité statistique, le caractère étudié et sa nature.
  \item Représenter cette distribution par un histogramme et tracer la courbe des effectifs cumulés croissants.
  \item Calculer l'âge moyen $\bar{x}$, la variance $V$ et l'écart-type $\sigma$.
  \item Déterminer la classe modale et le mode, puis la médiane par interpolation linéaire.
  \item Comparer moyenne, médiane et mode ; que peut-on dire de la symétrie de la distribution ?
\end{enumerate}

\section*{Exercice 2 -- Tirages et variable aléatoire (5 points)}
Une urne contient $4$ boules blanches et $6$ boules noires. On tire $3$ boules.
Soit $X$ le nombre de boules blanches obtenues.
\begin{enumerate}
  \item Les tirages sont \emph{simultanés} (sans remise).
  \begin{enumerate}[label=\alph*)]
    \item Déterminer la loi de probabilité de $X$.
    \item Calculer $E(X)$ et $V(X)$. Tracer la fonction de répartition de $X$.
  \end{enumerate}
  \item Les tirages sont \emph{successifs avec remise}. Reconnaître la loi de $X$, donner
  $E(X)$, $V(X)$ et $P(X \geqslant 1)$.
\end{enumerate}

\section*{Exercice 3 -- Probabilités conditionnelles (4 points)}
Dans un pays, $55\,\%$ des femmes en union vivent en milieu urbain. Parmi les urbaines,
$40\,\%$ utilisent une méthode contraceptive moderne ; parmi les rurales, $20\,\%$.
On choisit une femme au hasard. On note $U$ : « elle est urbaine » et $C$ : « elle utilise une méthode moderne ».
\begin{enumerate}
  \item Construire l'arbre pondéré de la situation.
  \item Calculer $P(C)$.
  \item La femme utilise une méthode moderne. Quelle est la probabilité qu'elle soit urbaine ?
  \item Les événements $U$ et $C$ sont-ils indépendants ? Justifier.
\end{enumerate}

\section*{Exercice 4 -- Couple de variables aléatoires (5 points)}
Soit $(X,Y)$ un couple de variables aléatoires dont la loi conjointe est :
\begin{center}
\begin{tabular}{c|ccc}
$X \backslash Y$ & 0 & 1 & 2\\ \hline
0 & 0{,}10 & 0{,}20 & 0{,}10\\
1 & 0{,}15 & 0{,}25 & 0{,}20
\end{tabular}
\end{center}
\begin{enumerate}
  \item Déterminer les lois marginales de $X$ et de $Y$.
  \item Calculer $E(X)$, $E(Y)$, $V(X)$, $V(Y)$.
  \item Calculer $E(XY)$ puis $\operatorname{Cov}(X,Y)$. $X$ et $Y$ sont-elles indépendantes ?
  \item Déterminer la loi de $S = X + Y$ et vérifier que $E(S) = E(X) + E(Y)$.
\end{enumerate}

\newpage
\begin{center}{\Large\bfseries Corrigé}\end{center}

\subsection*{Exercice 1}
\begin{enumerate}
  \item Population : les 100 femmes enquêtées ; unité : une femme ; caractère : âge à la première naissance, quantitatif continu (regroupé en classes).
  \item Classes de même amplitude (5 ans) : hauteurs proportionnelles aux effectifs. Effectifs cumulés : 12 ; 40 ; 75 ; 93 ; 100.
  \item Centres : 17,5 ; 22,5 ; 27,5 ; 32,5 ; 37,5.
  $\bar{x} = \dfrac{2650}{100} = 26{,}5$ ans ; $\overline{x^2} = 731{,}75$ ;
  $V = 731{,}75 - 26{,}5^2 = 29{,}5$ ; $\sigma \approx 5{,}43$ ans.
  \item Classe modale $[25;30[$ ; $M_o = 25 + 5\times\dfrac{7}{7+17} \approx 26{,}46$ ans.
  Médiane dans $[25;30[$ : $M_e = 25 + 5\times\dfrac{50-40}{35} \approx 26{,}43$ ans.
  \item $\bar{x} \approx M_e \approx M_o$ : distribution sensiblement symétrique.
\end{enumerate}

\subsection*{Exercice 2}
\begin{enumerate}
  \item a) $\binom{10}{3} = 120$ tirages.
  $P(X=0) = \frac{20}{120} = \frac16$, $P(X=1) = \frac{4\times15}{120} = \frac12$,
  $P(X=2) = \frac{6\times6}{120} = \frac{3}{10}$, $P(X=3) = \frac{4}{120} = \frac1{30}$.\\
  b) $E(X) = 1{,}2$ ; $E(X^2) = 2$ ; $V(X) = 2 - 1{,}44 = 0{,}56$.
  $F$ est en escalier : $0$ pour $x<0$ ; $\frac16$ sur $[0;1[$ ; $\frac23$ sur $[1;2[$ ; $\frac{29}{30}$ sur $[2;3[$ ; $1$ pour $x\geqslant3$.
  \item $X \sim \mathcal{B}(3\,;0{,}4)$ : $E(X) = 1{,}2$, $V(X) = 0{,}72$,
  $P(X\geqslant1) = 1 - 0{,}6^3 = 0{,}784$.
\end{enumerate}

\subsection*{Exercice 3}
\begin{enumerate}
  \setcounter{enumi}{1}
  \item $P(C) = 0{,}55\times0{,}40 + 0{,}45\times0{,}20 = 0{,}22 + 0{,}09 = 0{,}31$.
  \item $P_C(U) = \dfrac{0{,}22}{0{,}31} \approx 0{,}71$.
  \item $P(U\cap C) = 0{,}22 \neq P(U)P(C) = 0{,}1705$ : non indépendants.
\end{enumerate}

\subsection*{Exercice 4}
\begin{enumerate}
  \item $X$ : $P(X=0) = 0{,}4$, $P(X=1) = 0{,}6$. \quad
  $Y$ : $P(Y=0) = 0{,}25$, $P(Y=1) = 0{,}45$, $P(Y=2) = 0{,}30$.
  \item $E(X) = 0{,}6$ ; $V(X) = 0{,}24$ ; $E(Y) = 1{,}05$ ; $E(Y^2) = 1{,}65$ ; $V(Y) = 0{,}5475$.
  \item $E(XY) = 1\times0{,}25 + 2\times0{,}20 = 0{,}65$ ; $\operatorname{Cov}(X,Y) = 0{,}65 - 0{,}63 = 0{,}02$.
  Non indépendantes : $P(X=0,Y=1) = 0{,}20 \neq 0{,}4\times0{,}45 = 0{,}18$.
  \item $P(S=0) = 0{,}10$, $P(S=1) = 0{,}35$, $P(S=2) = 0{,}35$, $P(S=3) = 0{,}20$ ;
  $E(S) = 1{,}65 = 0{,}6 + 1{,}05$.
\end{enumerate}

\end{document}
